\documentclass[11pt,a4paper]{article}
\usepackage[T1]{fontenc}
\usepackage{lmodern}
\usepackage{amsmath,amssymb}
\usepackage{booktabs}
\usepackage{array}
\usepackage[a4paper,margin=2.5cm]{geometry}
\usepackage[hidelinks]{hyperref}
\usepackage{amsthm}

\newif\ifshowworked \showworkedtrue
\newenvironment{worked}{\ifshowworked\par\smallskip\noindent\textit{Working.}\par\else\setbox0\vbox\bgroup\fi}{\ifshowworked\par\smallskip\else\egroup\fi}

\newcommand{\Nc}{\mathcal{N}}
\newcommand{\Gc}{\mathcal{G}}
\newcommand{\eps}{\varepsilon}
\newcommand{\ii}{\mathrm{i}}
\newcommand{\bx}{\hat u_x}
\newcommand{\bz}{\hat u_z}
\newcommand{\bp}{\hat\varphi}
\newcommand{\bpsi}{\hat\psi}
\newcommand{\bt}{\hat\theta}
\newcommand{\kB}{k_{\mathrm{th}}}
\newcommand{\aME}{\alpha}

\title{\bfseries Lamb waves in a magneto-electro-elastic plate of the\\[2pt]
BaTiO$_3$/CoFe$_2$O$_4$ type with Moore--Gibson--Thompson heat\\[2pt]
conduction and the nonlocal strain gradient size effect}
\author{Dr.\ Vipin\\ Gurugram University}
\date{30 October 2026}

\begin{document}
\maketitle

\begin{abstract}
\noindent
A magneto-electro-elastic composite of barium titanate and cobalt ferrite carries three couplings in one
material: piezoelectric, piezomagnetic, and a magnetoelectric cross-coupling that neither phase has on its
own. We write the five coupled field equations of a plate of this material with poling along the normal,
reduce them to a plane-wave problem in the thickness coordinate, and carry the reduction to the boundary
matrix. The plate has thickness $2h$ and free faces; the wave is $\exp[\ii(kx-\omega t)]$ with complex $k$.
Heat conduction is of Moore--Gibson--Thompson type with pyroelectric and pyromagnetic terms; the elastic
stress carries the nonlocal strain gradient operator. The file gives the effective constants under two
homogenisation schemes, the constitutive components, the five equations, the $5\times5$ matrix, its degree
and parity, the amplitude ratios, the seven conditions per face, the boundary matrix and the hand limits,
including the thin-plate magnetoelectric coefficient.
\end{abstract}

\noindent\textit{Keywords.} magneto-electro-elastic; BaTiO$_3$/CoFe$_2$O$_4$; Lamb wave;
Moore--Gibson--Thompson; nonlocal strain gradient; magnetoelectric coefficient.

\section{Introduction}

A magneto-electro-elastic material responds to three inputs at once. Push it and it produces an electric
field and a magnetic field; apply an electric field and it strains; apply a magnetic field and it strains
too. The reason is composite construction: barium titanate is piezoelectric, cobalt ferrite is
magnetostrictive, and when the two are sintered together the strain of one phase is felt by the other, which
produces an electric response to a magnetic input and the other way round. This product property is called
the magnetoelectric effect, and it is the reason the composite is more interesting than either phase.

The effect is weak but real, and it is useful precisely because it is a cross-correlation: a magnetic field
can be measured through an electric voltage, without a coil. Lamb waves make it stronger, because a guided
wave concentrates energy in the plate and because the plate has a resonance structure that a bulk sample
does not. The price is that the theory becomes a five-field problem: two displacements, an electric
potential, a magnetic potential and a temperature.

Heat enters through three doors. The temperature changes the stress through the thermal expansion, it
changes the electric displacement through the pyroelectric constant, and it changes the magnetic induction
through the pyromagnetic constant. Both pyroelectric and pyromagnetic terms appear in plate papers on this
material, and we keep both \cite{Tocci2021}; the list of constants in that source includes the elastic,
piezoelectric, piezomagnetic, dielectric, magnetic, magnetoelectric, pyroelectric and pyromagnetic
coefficients in one place, which is what we need for the notation table. We read their Section~2 for the
list and we did not read a numbered equation from it.

Heat conduction itself is taken in the Moore--Gibson--Thompson form, with one time derivative cancelled. The
operator is read in \cite[Eqs.~(2.5) and (3.1)--(3.5)]{Quintanilla2019}. In a magneto-thermoelastic medium with memory
the same operator was used by Matlam and Kedar \cite{Matlam2026}; we read the title and the
abstract only, so no equation number is quoted from them.

The size effect is the nonlocal strain gradient theory, applied to the elastic part. Two lengths enter, a
gradient length $l_g$ that stiffens and a nonlocal length $l_{\rm nl}$ that softens. For a composite the
gradient length has a physical meaning that is easy to state: it is the distance over which the strain field
of one phase is felt by its neighbour, of the order of the phase size. That is also why the two
homogenisation schemes of Section~\ref{sec:A} matter: the effective constants depend on the volume fraction,
and the size lengths depend on the phase size.

Waves in magneto-electro-elastic plates are a well developed subject. Xiao, Han, Liu and Li \cite{Xiao2016}
study guided waves in magneto-electro-elastic plates; Kondaiah, Shankar and Ganesan \cite{Kondaiah2013}
study the same class with thermal effects; Kakar and Kakar \cite{Kakar2015} treat edge waves in an
electro-magneto-thermoelastic plate. We read titles and abstracts for all three and say so; no equation
number is quoted from them. The constitutive law we use is the standard linear one for this class, as
written by Chen, Pan and Chen \cite{Chen2007}; their paper was not opened, and the citation carries a note
to that effect. The reader who wants the numbered equations should go to \cite{Chen2007} Section~2 and
write them beside ours.

Our claim, in the words agreed with the supervisor, is this. Magnetoelectric coupling is kept as a full
constitutive constant with its own unit and sign, not as a perturbation. Heat conduction is of
Moore--Gibson--Thompson type with both the pyroelectric and the pyromagnetic terms, and the size effect is
kept in the gradient form on the elastic part only. Boundary conditions include the double tractions of the
gradient theory and both the electric and the magnetic conditions in open and closed versions. Parity is
the last part of the claim: this plate always splits into symmetric and antisymmetric families, unlike the
piezoelectric semiconductor plate of the companion paper, and Section~\ref{sec:D} explains why.

We also record the agreed downgrade. This paper is aimed at a second-quartile journal, not a first-quartile
one, and the decision between the two is to be taken at week nine on the ridge test, after the first
dispersion and attenuation curves exist. The ridge test asks whether the magnetoelectric coupling leaves a
visible signature on the attenuation ridge, not merely on the phase velocity.

\section{Setting and assumptions}

The plate occupies $-h\le z\le h$, is unbounded in $x$ and $y$, and has free faces. The material is a
two-phase BaTiO$_3$/CoFe$_2$O$_4$ composite, poled along $z$. Polling along the normal is what makes the
problem plane: the effective medium is transversely isotropic with the axis along $z$, class $6mm$ for the
elastic and electric parts, and the magnetic part has the same structure.

Plane strain in the $x$--$z$ plane: $u_y=0$, $\partial_y=0$. The wave is $\exp[\ii(kx-\omega t)]$ with real
$\omega$ and complex $k=k_r+\ii k_i$; positive $k_i$ means decay along $+x$. Across the thickness every
field is a sum of $\exp(\eta z)$.

The electromagnetic part is quasi-static: the wave speed is five orders of magnitude below the speed of
light, so the magnetic induction is the gradient of a scalar potential, $H_i=-\psi_{,i}$, and no wave
equation for the fields is solved. There are no free charges and no free currents in the plate, so both
Gauss laws are homogeneous. This is the standard magneto-electro-elastic approximation and it is a textbook
statement; no equation number is claimed.

The theory is linear. The temperature change $\theta$ is small, the fields are small, and the effective
constants are those of the homogenised medium, not of either phase. Recombination and carrier effects do not
exist here; the plate is an insulator and a dielectric.

The two size operators act as in the companion papers:
$\Nc(\eta)=1-l_{\rm nl}^2(\eta^2-k^2)$ on the inertia and $\Gc(\eta)=1-l_g^2(\eta^2-k^2)$ on the elastic
stress. The electric, magnetic and thermal parts carry no size operator.

\section{Notation}\label{sec:not}

\begin{center}\footnotesize
\begin{tabular}{p{2.0cm}p{5.6cm}p{1.9cm}p{4.6cm}}
\toprule
symbol & meaning in words & unit & where it enters\\
\midrule
$h$ & half the plate thickness & m & boundary rows\\
$k,\omega,\eta$ & wave number, angular frequency, thickness rate & rad/m, rad/s, 1/m & everywhere\\
$u_x,u_z$ & elastic displacements & m & rows 1, 2\\
$\varphi$ & electric potential; $\mathbf E=-\nabla\varphi$ & V & row 3\\
$\psi$ & magnetic potential; $\mathbf H=-\nabla\psi$ & A & row 4\\
$\theta$ & temperature change & K & row 5\\
$c_{11},c_{13},c_{33},c_{44}$ & effective elastic stiffnesses & Pa & rows 1, 2\\
$e_{31},e_{33},e_{15}$ & effective piezoelectric constants & C/m$^2$ & rows 1--3\\
$q_{31},q_{33},q_{15}$ & effective piezomagnetic constants & N/(A\,m) & rows 1, 2, 4\\
$\epsilon_{11},\epsilon_{33}$ & effective permittivities & F/m & row 3\\
$\mu_{11},\mu_{33}$ & effective magnetic permeabilities & H/m & row 4\\
$\aME_{11},\aME_{33}$ & effective magnetoelectric coefficients & s/m & rows 3, 4\\
$\beta_1,\beta_3$ & thermal stress coefficients & Pa/K & rows 1, 2, 5\\
$p_3$ & pyroelectric constant & C/(m$^2$K) & rows 3, 5\\
$\lambda_3$ & pyromagnetic constant & T/K & rows 4, 5\\
$\rho$ & density & kg/m$^3$ & rows 1, 2\\
$c_p$, $\kB$ & specific heat, thermal conductivity & J/(kg\,K), W/(m\,K) & row 5\\
$T_0$ & reference temperature & K & row 5\\
$\tau_q$ & MGT relaxation time & s & row 5\\
$l_g$, $l_{\rm nl}$ & size lengths of NSGT & m & $\Gc$, $\Nc$\\
$v$ & volume fraction of the ferrite phase & --- & effective constants only\\
$\mathbf M(\eta)$, $\mathbf B$ & amplitude and boundary matrices & --- & Sections~\ref{sec:C}, \ref{sec:G}\\
\bottomrule
\end{tabular}
\end{center}

Two units deserve a word because they are where mistakes happen in this literature. The magnetoelectric
coefficient $\aME$ appears in $\mathbf D=\dots+\aME\mathbf H$; with $\mathbf D$ in C/m$^2$ and $\mathbf H$ in
A/m, the unit of $\aME$ is s/m. Some papers write the coefficient in the magnetic equation instead, or as a
voltage coefficient with unit V\,m/A; before you compare numbers with any paper, convert to s/m and say which
convention you used. The piezomagnetic constant $q$ appears in $\boldsymbol\sigma=\dots-q\mathbf H$; with
stress in Pa and $\mathbf H$ in A/m its unit is Pa\,m/A, which is the same as N/(A\,m) or T.

\section{Source equations}\label{sec:src}

\subsection{Strain and fields}

$\eps_{ij}=\tfrac12(u_{i,j}+u_{j,i})$, $E_i=-\varphi_{,i}$, $H_i=-\psi_{,i}$. Under the plane wave,
$\eps_{xx}=\ii k\bx$, $\eps_{zz}=\eta\bz$, $\eps_{xz}=\tfrac12(\eta\bx+\ii k\bz)$, $E_x=-\ii k\bp$,
$E_z=-\eta\bp$, $H_x=-\ii k\bpsi$, $H_z=-\eta\bpsi$. Tag: textbook laws.

\subsection{Linear constitutive law of the magneto-electro-elastic medium}

\begin{equation}\label{eq:const}
\begin{aligned}
\sigma_{ij}&=c_{ijkl}\bigl(1-l_g^2\nabla^2\bigr)\eps_{kl}-e_{kij}E_k-q_{kij}H_k-\beta_{ij}\theta,\\
D_i&=e_{ikl}\eps_{kl}+\epsilon_{ik}E_k+\aME_{ik}H_k+p_i\theta,\\
B_i&=q_{ikl}\eps_{kl}+\aME_{ki}E_k+\mu_{ik}H_k+\lambda_i\theta .
\end{aligned}
\end{equation}
The structure is the standard one for this class, as written by Chen, Pan and Chen \cite{Chen2007}; we did
not open that paper, so no equation number is quoted, and the .bib entry carries a note. In \cite[Sec.~2]{Tocci2021} the list of
constants that survive, including the pyroelectric $p$ and the pyromagnetic $\lambda$, matches ours, which we did read; the coefficients are
$e,q,\zeta,p,\lambda$ in their notation, where $\zeta$ is the magnetoelectric coefficient. The gradient
operator on the elastic strain is the NSGT placement, ours, and the two thermal terms are textbook.
Tag: taken from the class literature without an equation number; the size operator is ours.

The magnetoelectric tensor is symmetric in the sense $\aME_{ik}=\aME_{ki}$, which is the Onsager symmetry of
the pair: the coefficient that makes a magnetic field produce electric displacement is the same one that
makes an electric field produce magnetic induction. This is why one symbol appears in two places in
\eqref{eq:const}, and it is the first thing to check if a computed dispersion curve looks asymmetric.

\subsection{Momentum, Gauss laws and heat}

\begin{equation}\label{eq:gov}
\bigl(1-l_{\rm nl}^2\nabla^2\bigr)\rho\ddot u_i=\sigma_{ij,j},\qquad
D_{i,i}=0,\qquad B_{i,i}=0,
\end{equation}
\begin{equation}\label{eq:heat}
\Bigl(1+\tau_q\frac{\partial}{\partial t}\Bigr)
\Bigl[\rho c_p\dot\theta+T_0\beta_{ij}\dot\eps_{ij}+T_0p_i\dot E_i+T_0\lambda_i\dot H_i\Bigr]
=\kB\nabla^2\theta .
\end{equation}
The momentum balance with the nonlocal operator is taken from \cite[Eqs.~(18)--(19)]{Barretta2018}, who
attribute the differential law to Lim, Zhang and Reddy \cite{Lim2015}; we did not open \cite{Lim2015} and say
so in the .bib. Both Gauss laws are textbook electrostatics and magnetostatics of a dielectric with no
free charge. The heat equation is the Moore--Gibson--Thompson equation of
\cite[Eqs.~(2.5) and (3.1)--(3.5)]{Quintanilla2019} with the two caloric terms added; the pyroelectric term is
classical and the pyromagnetic one is its magnetic partner, and both are the reciprocity partners of the
$+p_3\theta$ and $+\lambda_3\theta$ terms in \eqref{eq:const}.
Tag: the momentum law is taken from \cite{Barretta2018}; the heat operator from \cite{Quintanilla2019}; the
placement of the operators and the pairing of the caloric terms are ours.

\section{Calculation A: effective constants and constitutive components}\label{sec:A}

\subsection{Two homogenisation schemes}

The composite is described by effective constants. We fix two schemes, and we state them so that you can
reproduce both and see how much the dispersion curve depends on the choice.

Scheme one is the rule of mixtures. Every constant is the volume average of the two phases,
\begin{equation}\label{eq:voigt}
\bar C=v\,C^{(f)}+(1-v)\,C^{(m)},\qquad
\bar e=v\,e^{(f)},\qquad \bar q=v\,q^{(f)},\qquad
\bar\eps=v\,\eps^{(f)}+(1-v)\,\eps^{(m)},\qquad \bar\mu=v\,\mu^{(f)}+(1-v)\,\mu^{(m)},
\end{equation}
where the superscript $(f)$ is the ferrite, which carries the piezomagnetic constants, and $(m)$ is the
barium titanate, which carries the piezoelectric ones. This is a textbook mixture rule and needs no
citation. It has a defect that matters here: both phases have zero magnetoelectric coefficient, so the
average \eqref{eq:voigt} gives $\bar\aME=0$ and the very effect we want to study disappears.

Scheme two is a two-phase estimate of the Mori--Tanaka type. Write the strain concentration tensor $\mathbb A$
of the inclusion phase, and use one weight for all the coupled constants,
\begin{equation}\label{eq:mt}
\bar{\mathbb C}={\mathbb C}^{(m)}+v\bigl({\mathbb C}^{(f)}-{\mathbb C}^{(m)}\bigr)\mathbb A,\qquad
\bar e={\mathbb C}^{(m)}\mathbb A\,{\mathbb e}^{(m)}\ \text{replaced by}\ v\,{\mathbb e}^{(f)}\mathbb A,
\end{equation}
with the same weight applied to $q$, $\eps$ and $\mu$. For the magnetoelectric coefficient we take the
product form
\begin{equation}\label{eq:alpha}
\bar\aME_{33}=v(1-v)\,\frac{e^{(f)}_{31}q^{(m)}_{33}}{\bar C_0},
\end{equation}
where $\bar C_0$ is an effective stiffness of the composite used as the normaliser. The units check: $\bar
e$ is C/m$^2$, $\bar q$ is N/(A\,m), and their product divided by a stiffness in Pa gives C/(A\,m), which is
s/m, the unit of $\aME$ fixed in the notation table.

We must be plain about the status of \eqref{eq:mt} and \eqref{eq:alpha}. The product form is the standard
reason why a two-phase composite shows a magnetoelectric effect that neither phase shows, but we did not read
a numbered equation for it in any source, and the normaliser $\bar C_0$ is not fixed by anything we read.
Before submission, find the formula in \cite[Sec.~2]{Chen2007} or in \cite[Sec.~2]{Tocci2021}, quote its
number, and replace \eqref{eq:alpha} by the published form. Until then, treat $\bar\aME_{33}$ as a parameter
of the model and say so.

The sign of $\aME$ is a convention, and it is the convention that trips people. Define the positive sense of
$z$ as the poling direction of the piezoelectric phase. Reversing that poling reverses $e_{31}$, and by
\eqref{eq:alpha} it reverses $\aME_{33}$; the physics does not change, the sign of the predicted induced
field does. Write the convention in the caption of every figure.

\subsection{The two schemes, element by element}

Scheme one is short enough to write out for the constants we need, and writing it out shows what volume
fraction does to each of them.

\begin{worked}
Let $v$ be the ferrite fraction. The elastic stiffnesses of the composite are
$\bar c_{11}=v c^{(f)}_{11}+(1-v)c^{(m)}_{11}$ and the same for $c_{13},c_{33},c_{44}$; the piezoelectric
constants come only from the barium titanate, $\bar e_{31}=(1-v)e^{(m)}_{31}$ and the same for $e_{33},e_{15}$;
the piezomagnetic constants come only from the ferrite, $\bar q_{31}=vq^{(f)}_{31}$ and the same for
$q_{33},q_{15}$; the permittivity and the permeability are averages,
$\bar\epsilon_{33}=v\epsilon^{(f)}_{33}+(1-v)\epsilon^{(m)}_{33}$ and
$\bar\mu_{33}=v\mu^{(f)}_{33}+(1-v)\mu^{(m)}_{33}$; the two pyroelectric and pyromagnetic coefficients follow
their phases: $\bar p_3=(1-v)p^{(m)}_3$ and $\bar\lambda_3=v\lambda^{(f)}_3$. The magnetoelectric coefficient
is zero by construction, because both phases have zero and an average of zeros is zero.
\end{worked}

\noindent Read the consequences off these lines before you go to scheme two. The piezoelectric constant falls
with the ferrite fraction, the piezomagnetic one rises with it, so their product is largest in the middle,
and the magnetoelectric effect of the composite vanishes at both ends. That is the physics that the
magnetoelectric coefficient must carry, and scheme two is the smallest change to scheme one that gives it a
non-zero value.

Scheme two replaces the arithmetic average by a weighted one. Write $\mathbb A$ for the strain concentration
tensor of the inclusion phase, which is a fourth-order tensor built from the stiffness of the matrix and the
shape of the inclusion. Every constant is then weighted by $\mathbb A$ rather than by the volume fractions
alone, and the magnetoelectric coefficient takes the product form of \eqref{eq:alpha}. The shape of the
inclusion is what makes $\mathbb A$ different from the identity, so the difference between the two schemes is
largest for elongated phases, which is also where the measured magnetoelectric coefficient of such composites
is largest.

\subsection{Constitutive components in our geometry}

\begin{worked}
Apply the routine: $\partial_x\to\ii k$, $\partial_z\to\eta$, Laplacian $\to\eta^2-k^2$, and $-e_{kij}E_k$
with $E_z=-\eta\bp$ gives $+\eta e_{33}\bp$ for the $zz$ component and $-e_{31}E_z=+\eta e_{31}\bp$ for the
$xx$ component. The piezomagnetic terms follow the same pattern with $q$ in place of $e$ and $\bpsi$ in place
of $\bp$. In $D_x$ the electric term is $-\ii k\epsilon_{11}\bp$ and the magnetoelectric term is
$-\ii k\aME_{11}\bpsi$; in $D_z$ they are $-\eta\epsilon_{33}\bp$ and $-\eta\aME_{33}\bpsi$. The magnetic
components are the same with $D\to B$, $\bp\to\bpsi$, and the pair $(\epsilon,\aME)$ replaced by
$(\aME,\mu)$.
\end{worked}

\noindent
\begin{align}
\hat\sigma_{xx}&=\Gc\bigl(\ii kc_{11}\bx+\eta c_{13}\bz\bigr)+\eta e_{31}\bp+\eta q_{31}\bpsi-\beta_1\bt,
\label{eq:sxx}\\
\hat\sigma_{zz}&=\Gc\bigl(\ii kc_{13}\bx+\eta c_{33}\bz\bigr)+\eta e_{33}\bp+\eta q_{33}\bpsi-\beta_3\bt,
\label{eq:szz}\\
\hat\sigma_{xz}&=\Gc c_{44}\bigl(\eta\bx+\ii k\bz\bigr)+\ii ke_{15}\bp+\ii kq_{15}\bpsi,
\label{eq:sxz}\\
\hat D_x&=e_{15}\bigl(\eta\bx+\ii k\bz\bigr)-\ii k\epsilon_{11}\bp-\ii k\aME_{11}\bpsi,
\label{eq:Dx}\\
\hat D_z&=\ii ke_{31}\bx+\eta e_{33}\bz-\eta\epsilon_{33}\bp-\eta\aME_{33}\bpsi+p_3\bt,
\label{eq:Dz}\\
\hat B_x&=q_{15}\bigl(\eta\bx+\ii k\bz\bigr)-\ii k\aME_{11}\bp-\ii k\mu_{11}\bpsi,
\label{eq:Bx}\\
\hat B_z&=\ii kq_{31}\bx+\eta q_{33}\bz-\eta\aME_{33}\bp-\eta\mu_{33}\bpsi+\lambda_3\bt .
\label{eq:Bz}
\end{align}

Read the symmetry off \eqref{eq:Dz} and \eqref{eq:Bz}. The coefficient $\aME_{33}$ sits in the same place in
both, which is the Onsager pair; if you change its sign in one line and not the other, the model loses its
reciprocity and the parity argument of Section~\ref{sec:D} still holds, but the energy argument does not.
That is the kind of error that survives every numerical check and fails on the first referee report.

\section{Calculation B: the five governing equations}\label{sec:B}

\begin{worked}
The two momentum equations are the same routine as in the companion papers:
$-\bigl(\ii k\hat\sigma_{xx}+\eta\hat\sigma_{xz}\bigr)-\Nc\rho\omega^2\bx=0$ and
$-\bigl(\ii k\hat\sigma_{xz}+\eta\hat\sigma_{zz}\bigr)-\Nc\rho\omega^2\bz=0$. The two Gauss laws are
$\ii k\hat D_x+\eta\hat D_z=0$ and $\ii k\hat B_x+\eta\hat B_z=0$. The heat equation is
$f\bigl[-\ii\omega\rho c_p\bt+T_0\omega k\beta_1\bx-\ii\omega T_0\beta_3\eta\bz
+\ii\omega T_0p_3\eta\bp+\ii\omega T_0\lambda_3\eta\bpsi\bigr]-\kB L\bt=0$, obtained by the same
substitution: each time derivative is $-\ii\omega$, $T_0p_i\dot E_i$ gives
$T_0p_3(-\ii\omega)(-\eta\bp)=\ii\omega T_0p_3\eta\bp$, and the magnetic term gives the same with
$\lambda_3$ and $\bpsi$.
\end{worked}

\noindent
\begin{equation}\label{eq:g}
\begin{gathered}
-\bigl(\ii k\hat\sigma_{xx}+\eta\hat\sigma_{xz}\bigr)-\Nc\rho\omega^2\bx=0,\qquad
-\bigl(\ii k\hat\sigma_{xz}+\eta\hat\sigma_{zz}\bigr)-\Nc\rho\omega^2\bz=0,\\
\ii k\hat D_x+\eta\hat D_z=0,\qquad \ii k\hat B_x+\eta\hat B_z=0,\\
f\Bigl[-\ii\omega\rho c_p\bt+T_0\omega k\beta_1\bx-\ii\omega T_0\beta_3\eta\bz
+\ii\omega T_0p_3\eta\bp+\ii\omega T_0\lambda_3\eta\bpsi\Bigr]-\kB L\bt=0 .
\end{gathered}
\end{equation}

\section{Calculation C: the amplitude matrix}\label{sec:C}

Order the unknowns as $\mathbf a=(\bx,\bz,\bp,\bpsi,\bt)^{\top}$. Rows: momentum $x$, momentum $z$,
electric Gauss, magnetic Gauss, heat.

\subsection{Rows 1 and 2}

\begin{worked}
Put \eqref{eq:sxx} and \eqref{eq:sxz} into the first equation of \eqref{eq:g}. For $\bx$:
$\ii k\Gc c_{11}\ii k=-k^2\Gc c_{11}$ and $\eta\Gc c_{44}\eta=\eta^2\Gc c_{44}$, with the overall minus and
the inertia. For $\bz$: $\ii k\Gc c_{13}\eta+\eta\Gc c_{44}\ii k$, with a minus. For $\bp$:
$\ii k\eta e_{31}+\eta\ii ke_{15}$, with a minus. For $\bpsi$: the same with $q$. For $\bt$:
$\ii k(-\beta_1)$ with the overall minus gives $+\ii k\beta_1$. The second row is the same routine with
$c_{11}\leftrightarrow c_{44}$, $c_{13}\to c_{33}$, $e_{31}\to e_{33}$, $q_{31}\to q_{33}$ and the pairing
of the shear components.
\end{worked}

\noindent
\begin{equation}\label{eq:row12}
\begin{aligned}
&M_{11}=-\Nc\rho\omega^2+\Gc(c_{11}k^2-c_{44}\eta^2),\quad
M_{12}=-\ii k\eta\Gc(c_{13}+c_{44}),\\
&M_{13}=-\ii k\eta(e_{31}+e_{15}),\quad
M_{14}=-\ii k\eta(q_{31}+q_{15}),\quad M_{15}=\ii k\beta_1,\\
&M_{21}=-\ii k\eta\Gc(c_{13}+c_{44}),\quad
M_{22}=-\Nc\rho\omega^2+\Gc(c_{44}k^2-c_{33}\eta^2),\\
&M_{23}=e_{15}k^2-e_{33}\eta^2,\quad
M_{24}=q_{15}k^2-q_{33}\eta^2,\quad M_{25}=\eta\beta_3 .
\end{aligned}
\end{equation}

\subsection{Rows 3 and 4}

\begin{worked}
For the electric Gauss law, $\ii k\hat D_x+\eta\hat D_z$ with \eqref{eq:Dx} and \eqref{eq:Dz}. For $\bx$:
$\ii ke_{15}\eta+\eta\ii ke_{31}$. For $\bz$: $-k^2e_{15}+\eta^2e_{33}$. For $\bp$:
$\epsilon_{11}k^2-\epsilon_{33}\eta^2$. For $\bpsi$: $\aME_{11}k^2-\aME_{33}\eta^2$. For $\bt$: $\eta p_3$.
The magnetic row is the same with $(e,q)\to(q,e)$ care: the $\bx$ entry is $\ii kq_{15}\eta+\eta\ii kq_{31}$,
the $\bz$ entry is $-k^2q_{15}+\eta^2q_{33}$, the $\bp$ entry is $\aME_{11}k^2-\aME_{33}\eta^2$ again, the
$\bpsi$ entry is $\mu_{11}k^2-\mu_{33}\eta^2$, and the thermal entry is $\eta\lambda_3$.
\end{worked}

\noindent
\begin{equation}\label{eq:row34}
\begin{aligned}
&M_{31}=\ii k\eta(e_{15}+e_{31}),\quad M_{32}=e_{33}\eta^2-e_{15}k^2,\quad
M_{33}=\epsilon_{11}k^2-\epsilon_{33}\eta^2,\\
&M_{34}=\aME_{11}k^2-\aME_{33}\eta^2,\quad M_{35}=\eta p_3,\\
&M_{41}=\ii k\eta(q_{15}+q_{31}),\quad M_{42}=q_{33}\eta^2-q_{15}k^2,\quad
M_{43}=\aME_{11}k^2-\aME_{33}\eta^2,\\
&M_{44}=\mu_{11}k^2-\mu_{33}\eta^2,\quad M_{45}=\eta\lambda_3 .
\end{aligned}
\end{equation}

\subsection{Row 5}

\begin{worked}
Read the coefficients from the last line of \eqref{eq:g}: $M_{51}=fT_0\omega k\beta_1$,
$M_{52}=-\ii\omega fT_0\beta_3\eta$, $M_{53}=\ii\omega fT_0p_3\eta$, $M_{54}=\ii\omega fT_0\lambda_3\eta$,
$M_{55}=-\ii\omega f\rho c_p-\kB L$.
\end{worked}

\noindent
\begin{equation}\label{eq:row5}
M_{51}=fT_0\omega k\beta_1,\quad M_{52}=-\ii\omega fT_0\beta_3\eta,\quad
M_{53}=\ii\omega fT_0p_3\eta,\quad M_{54}=\ii\omega fT_0\lambda_3\eta,\quad
M_{55}=-\ii\omega f\rho c_p-\kB L .
\end{equation}

\noindent Collecting \eqref{eq:row12}--\eqref{eq:row5} gives $\mathbf M(\eta)\mathbf a=\mathbf 0$ with
$\Nc=1-l_{\rm nl}^2L$, $\Gc=1-l_g^2L$, $L=\eta^2-k^2$ and $f=1-\ii\omega\tau_q$. The matrix is not written
as one array because its entries are long; the five row equations are the matrix, in order.

The magnetic and electric rows are twins. Exchange $\bp\leftrightarrow\bpsi$, $e\leftrightarrow q$ and
$\epsilon\leftrightarrow\mu$, and one becomes the other, with $\aME$ unchanged. That symmetry is a useful
check: if the determinant of your $5\times5$ matrix is not invariant under that exchange when
$\epsilon=\mu$ and $e=q$, there is a sign error in row 3 or row 4.

\section{Calculation D: degree and parity}\label{sec:D}

\subsection{Degree}

\begin{worked}
The highest power of $\eta$ in each entry: row 1 has $M_{11}\sim\eta^4$ (when $l_g\ne0$), $M_{12}\sim\eta^3$,
$M_{13},M_{14}\sim\eta^1$, $M_{15}\sim\eta^0$; row 2 has $M_{21}\sim\eta^3$, $M_{22}\sim\eta^4$,
$M_{23},M_{24}\sim\eta^2$, $M_{25}\sim\eta^1$; rows 3 and 4 have entries of degree 1, 2, 2, 2, 1; row 5 has
degrees 0, 1, 1, 1, 2. The permutation with the largest total is the diagonal one, giving 14. No other
permutation can exceed it, because rows 1 and 2 give at most four each and rows 3, 4, 5 at most two each.
\end{worked}

\noindent
\begin{equation}\label{eq:deg}
\deg_\eta\det\mathbf M=\begin{cases}14,&l_g\ne0,\\10,&l_g=0,\end{cases}
\end{equation}
so the plate carries fourteen partial waves and needs seven conditions per face, or ten and five when the
gradient length is zero.

\subsection{Parity}

\begin{worked}
From row 1: $M_{12}$ odd gives $s_2=-s_1$; $M_{13}$ odd gives $s_3=-s_1$; $M_{14}$ odd gives $s_4=-s_1$;
$M_{15}$ even gives $s_5=s_1$. Take $s_1=-1$, so $s=(-1,+1,+1,+1,-1)$. Check every other entry: row 2 has
$M_{21}$ odd, needing $s_2s_1=-1$; $M_{23}$ and $M_{24}$ even, needing $s_2s_3=s_2s_4=+1$; $M_{25}$ odd,
needing $s_2s_5=-1$; all these hold. Row 3 has $M_{31}$ odd, $M_{32}$, $M_{33}$, $M_{34}$ even, $M_{35}$ odd,
and with $s_1=-1$, $s_5=-1$ and $s_2=s_3=s_4=+1$ every one of them holds. Row 4 is the twin of row 3, so it
holds too. Row 5 has $M_{51}$ even, $M_{52},M_{53},M_{54}$ odd, $M_{55}$ even; with $s_5=-1$ and
$s_2=s_3=s_4=+1$ each product is $-1$ where the entry is odd and $+1$ where it is even, so it holds.
\end{worked}

\noindent
\begin{equation}\label{eq:parity}
\det\mathbf M(-\eta)=\det\mathbf M(\eta)\quad\text{for every value of every coupling,}
\end{equation}
with the sign vector $(-1,+1,+1,+1,-1)$: the displacement along the plate and the temperature are odd
through the thickness; the displacement across it, the electric potential and the magnetic potential are
even.

Compare this with the companion paper on the piezoelectric semiconductor. There the carrier density sits in
the matrix in a place that no sign vector can accommodate together with the temperature, so the parity fails
as soon as the thermodiffusion or the Peltier coefficient is switched on. Here there is no carrier field: all
five couplings are of even parity in the appropriate sense, and the split survives. The magnetoelectric
coefficient, which one might suspect of breaking it, does not, because it appears only in the two Gauss rows
and always as an even entry.

The practical consequence is the factorisation of the boundary matrix,
\begin{equation}\label{eq:factor}
\det\mathbf B=\det\mathbf B_{\rm S}\cdot\det\mathbf B_{\rm A},
\end{equation}
with seven rows and seven columns in each block. Work with the two blocks separately; the full $14\times14$
determinant is then only a check.

\section{Calculation E: amplitude ratios and partial waves}\label{sec:E}

\begin{worked}
At a root $\eta_m$ the adjugate $\mathbf A$ of $\mathbf M(\eta_m)$ has columns proportional to the null
vector. Use the cofactor formula $A_{ij}=(-1)^{i+j}\det\mathbf M^{(ji)}$, where $\mathbf M^{(ji)}$ is the
$4\times4$ matrix obtained by deleting row $j$ and column $i$. Take the column with the largest entries and
divide by its largest entry. This avoids dividing by a small determinant and stays finite when two roots are
close.
\end{worked}

\noindent
\begin{equation}\label{eq:general}
\bigl(u_x,u_z,\varphi,\psi,\theta\bigr)=\sum_{m=1}^{N}a_m
\bigl(\rho^{(1)}_m,\rho^{(2)}_m,\rho^{(3)}_m,\rho^{(4)}_m,1\bigr)e^{\eta_mz}e^{\ii(kx-\omega t)},
\end{equation}
with $N=14$ or $10$ as in \eqref{eq:deg}. Because the parity \eqref{eq:parity} holds, the roots come in pairs
and the amplitude ratios of a pair differ only by the sign vector: the odd fields flip, the even ones do not.
Each pair can therefore be written as a hyperbolic cosine for the even fields and a hyperbolic sine for the
odd ones, which is the classical form of the symmetric and antisymmetric solutions.

\section{Calculation F: boundary conditions}\label{sec:F}

At each face, with outward normal $\mathbf n=(0,0,\pm1)$:

\begin{enumerate}
\item[(b1)] normal traction with the surface gradient of the double traction:
$\bigl[\sigma_{zz}-\partial_xR_{xz}\bigr]_{z=z_0}=0$;
\item[(b2)] shear traction with its surface gradient:
$\bigl[\sigma_{xz}-\partial_xR_{xx}\bigr]_{z=z_0}=0$;
\item[(b3)] normal double traction: $\bigl[R_{zz}\bigr]_{z=z_0}=0$;
\item[(b4)] in-plane double traction: $\bigl[R_{zx}\bigr]_{z=z_0}=0$;
\item[(b5)] electric condition: open circuit $\bigl[D_z\bigr]_{z=z_0}=0$, or closed circuit
$\bigl[\varphi\bigr]_{z=z_0}=0$;
\item[(b6)] magnetic condition: open circuit $\bigl[B_z\bigr]_{z=z_0}=0$, or closed circuit
$\bigl[\psi\bigr]_{z=z_0}=0$;
\item[(b7)] thermal condition: adiabatic $\bigl[\partial_z\theta\bigr]_{z=z_0}=0$, or isothermal
$\bigl[\theta\bigr]_{z=z_0}=0$.
\end{enumerate}

The surface-gradient form of (b1) and (b2) and the double-traction conditions (b3) and (b4) are taken from
\cite[Sec.~2, Eq.~(14)]{Charalambopoulos2019} and extended by us to the electric, magnetic and thermal parts.
For the four electric and magnetic conditions we follow standard practice in this literature; titles and
abstracts of \cite{Xiao2016} and \cite{Kondaiah2013} were read, their equations were not, so we quote no
number. Open circuit on the magnetic side is the condition that matters for a magnetoelectric measurement,
because it is the condition under which the induced field is not shorted out by the surroundings. That
statement is ours, not a citation.

\begin{worked}
The nonlocal operator is inverted root by root: at root $\eta_m$ the constitutive law reads
$\Nc(\eta_m)\hat\sigma_m=\hat\sigma^{\rm loc}_m$, so the stress in a boundary row is
$\hat\sigma^{\rm loc}_m/\Nc(\eta_m)$ and the factor belongs to column $m$. No common operator can be pulled
out of a sum over roots.
\end{worked}

\noindent The seven row functions at $z=+h$, with the amplitudes of partial wave $m$ and
$\partial_x\to\ii k$, are
\begin{align}
b_1&=\frac{1}{\Nc(\eta_m)}\Bigl[\hat\sigma_{zz}-\ii k\,l_g^2c_{44}L\hat\eps_{xz}\Bigr],\quad
b_2=\frac{1}{\Nc(\eta_m)}\Bigl[\hat\sigma_{xz}-\ii k\,l_g^2c_{11}L\hat\eps_{xx}\Bigr],
\label{eq:b12}\\
b_3&=\frac{1}{\Nc(\eta_m)}l_g^2c_{33}L\hat\eps_{zz},\quad
b_4=\frac{1}{\Nc(\eta_m)}l_g^2c_{44}L\hat\eps_{xz},\quad
b_5=\hat D_z,\quad b_6=\hat B_z,\quad b_7=\eta_m\bt .
\label{eq:b34567}
\end{align}

\begin{worked}
At $z=-h$ the partial wave carries $\exp(-\eta_mh)$ instead of $\exp(\eta_mh)$, and the outward normal is
$(0,0,-1)$, so $\sigma_{zz}$, $\sigma_{xz}$, $D_z$ and $B_z$ change sign while the double tractions and the
temperature gradient behave as in the companion papers: the double tractions keep their sign because they are
quadratic in the normal, and the temperature gradient changes sign because $\partial_z$ is odd.
\end{worked}

\noindent With $c_m=\exp(\eta_mh)$ and $d_m=\exp(-\eta_mh)$,
\begin{equation}\label{eq:rows}
b_1c_m,\ b_2c_m,\ b_3c_m,\ b_4c_m,\ b_5c_m,\ b_6c_m,\ b_7c_m\ (z=+h),\qquad
-b_1d_m,\ -b_2d_m,\ b_3d_m,\ b_4d_m,\ -b_5d_m,\ -b_6d_m,\ -b_7d_m\ (z=-h).
\end{equation}

If you choose closed circuit for the electric condition, replace (b5) by $\bp$ alone, and the same for the
magnetic condition. The choice changes the dispersion curves much more than the size lengths do, so state it
in every figure caption.

\section{Calculation G: the boundary matrix}\label{sec:G}

\noindent Number the roots $\eta_1,\dots,\eta_N$, put the row functions \eqref{eq:b12} and \eqref{eq:b34567}
in the rows of \eqref{eq:rows}, and the roots in the columns. The dispersion relation is
\begin{equation}\label{eq:disp}
\det\mathbf B(k,\omega)=0,\qquad \mathbf B \text{ is }14\times14\ (l_g\ne0),\ 10\times10\ (l_g=0).
\end{equation}
When $l_g=0$ the rows (b3) and (b4) vanish and are deleted with the four extra roots. The parity
\eqref{eq:parity} then blocks the matrix into two $7\times7$ (or $5\times5$) problems as in
\eqref{eq:factor}.

\begin{worked}
The frequency is real and given; $k$ is complex. The dispersion relation is one complex equation in two real
unknowns, so its solutions are isolated pairs, and the search is in the complex $k$ plane started from the
uncoupled elastic solution. Multiply column $m$ by $\exp(-|\mathrm{Re}\,\eta_m|h)$ before evaluating any
determinant; the factor is common to the column, so it moves no zero, but it keeps the columns comparable
once the gradient roots appear.
\end{worked}

\noindent Two numerical cautions belong in the record. First, the magnetic and electric rows are nearly
proportional when $\aME$ is small, so the boundary matrix is nearly singular in a way that has nothing to do
with a mode; watch the smallest singular value against $k$, not the determinant alone. Second, the
piezomagnetic constants are large compared with the piezoelectric ones, so the magnetic block dominates the
scaling of the rows; normalise the rows before you look at any singular value.

\section{Calculation H: hand limits}\label{sec:H}

\subsection{Both couplings off}

\begin{worked}
Set $e=q=\aME=p_3=\lambda_3=0$ and $\beta_1=\beta_3=0$. Rows 3, 4 and 5 decouple from rows 1 and 2, and the
mechanical part is the $2\times2$ block of the companion papers, whose determinant gives the classical
Rayleigh--Lamb condition. With $c_{11}=c_{33}=\lambda+2\mu$, $c_{13}=\lambda$, $c_{44}=\mu$ the two families
obey
\end{worked}

\begin{equation}\label{eq:RL}
\frac{\tan(qh)}{\tan(ph)}=-\frac{4k^2pq}{(k^2-q^2)^2}\ ({\rm symmetric}),\qquad
\frac{\tan(qh)}{\tan(ph)}=-\frac{(k^2-q^2)^2}{4k^2pq}\ ({\rm antisymmetric}).
\end{equation}

\subsection{Electric off, magnetic on}

\begin{worked}
Set $e=\aME=p_3=0$ and keep $q$ and $\mu$. The electric row decouples entirely, since $D_z$ then carries only
$\bp$ and $\bpsi$ through $\aME=0$; the remaining $4\times4$ system is the piezomagnetic analogue of a
piezoelectric plate. Its two size-free limits are the same as in \eqref{eq:RL} with the magnetic stiffening
of the appropriate modulus in place of the electric one, and the check is that the magnetic open-circuit
speed is above the closed-circuit one, exactly as the electric open-circuit speed is above its
closed-circuit value.
\end{worked}

\subsection{Thin plate: the effective coupling}

\begin{worked}
Take the frequency--thickness product small, so the fields are almost uniform across the thickness, and keep
all couplings. The faces are free and open on both the electric and the magnetic side, so
$\sigma_{zz}=D_z=B_z=0$ and the plate carries only $\eps_{xx}$. From \eqref{eq:Dz} and \eqref{eq:Bz} with
$\theta=0$,
\[
\epsilon_{33}E_z+\aME_{33}H_z=-e_{31}\eps_{xx},\qquad
\aME_{33}E_z+\mu_{33}H_z=-q_{31}\eps_{xx},
\]
Solve the $2\times2$ system:
\[
E_z=\frac{-e_{31}\mu_{33}+\aME_{33}q_{31}}{\epsilon_{33}\mu_{33}-\aME_{33}^2}\,\eps_{xx},\qquad
H_z=\frac{-q_{31}\epsilon_{33}+\aME_{33}e_{31}}{\epsilon_{33}\mu_{33}-\aME_{33}^2}\,\eps_{xx}.
\]
Put both into $\sigma_{xx}$ of \eqref{eq:sxx} with $\eps_{zz}=0$ and $\theta=0$:
\[
\bar c_{11}=c_{11}+\frac{e_{31}^2\mu_{33}-2e_{31}q_{31}\aME_{33}+q_{31}^2\epsilon_{33}}
{\epsilon_{33}\mu_{33}-\aME_{33}^2}.
\]
\end{worked}

\noindent The thin-plate modulus is the piezoelectric one plus a magnetic correction and a magnetoelectric
cross term, and the cross term enters with a minus sign because it is the product of the two couplings
through the same $\aME$. The ratio of the two induced fields in the same open-circuit state is
\begin{equation}\label{eq:ratio}
\frac{H_z}{E_z}=\frac{\aME_{33}e_{31}-q_{31}\epsilon_{33}}{\aME_{33}q_{31}-e_{31}\mu_{33}},
\end{equation}
which is the thin-plate magnetoelectric response. Its unit is A/V, that is siemens, and it is large when the
two couplings are large and the determinant $\epsilon_{33}\mu_{33}-\aME_{33}^2$ is small. This is the
quantity to plot against the volume fraction: it must vanish at $v=0$ and at $v=1$ because both phases alone
have no magnetoelectric effect, so it has a maximum somewhere in between, and that maximum is the physical
content of the composite.

\subsection{The two size lengths and the MGT time}

\begin{worked}
Setting $l_g=l_{\rm nl}=0$ returns the local plate with ten roots; setting $\tau_q=0$ returns the Fourier
heat equation. Increasing $l_g$ stiffens the elastic terms of rows 1 and 2, increasing $l_{\rm nl}$ softens
the inertia. The two effects are opposite and are separated by their order in $k$: the gradient term grows
with the square of the wave number, the nonlocal term with the square of the frequency, so a dispersion curve
at fixed frequency separates them.
\end{worked}

\noindent The MGT time constant is set by the thickness, $\tau_T=(2h)^2/(\pi^2\kappa)$ with $\kappa=\kB/(\rho
c_p)$ the thermal diffusivity. Its effect on the attenuation is a shift of the classical thermoelastic peak
and a change of its height; the shift is the measurable signature and it is the reason the operator is worth
keeping in a paper that is otherwise about the magnetoelectric effect.

\section{Self-test list}\label{sec:tests}

\begin{enumerate}
\item Units inside each row of the matrix, in particular the magnetoelectric entries, which must combine
with the dielectric ones to give the same unit as $D_z$ per unit potential.
\item Degree: fourteen with $l_g\ne0$, ten with $l_g=0$.
\item Parity: the sign vector $(-1,+1,+1,+1,-1)$ for every choice of couplings, and the factorisation of the
assembled matrix.
\item The electric--magnetic exchange symmetry of rows 3 and 4 when $\epsilon=\mu$ and $e=q$.
\item The classical limit \eqref{eq:RL} with all couplings off.
\item The open-circuit ordering, electric and magnetic, against the closed-circuit values.
\item The thin-plate modulus and the field ratio of \eqref{eq:ratio}, both vanishing with $\aME$.
\item The volume-fraction limits: the effective constants must return the pure phases at $v=0$ and $v=1$,
and the magnetoelectric coefficient must vanish at both ends.
\item The two size lengths in opposite directions, and the MGT time shifting the thermal peak.
\item The boundary switches: closed circuit electric, closed circuit magnetic, isothermal face, one at a
time, each changing one row.
\end{enumerate}

\section{Work plan}

\begin{enumerate}
\item Week 1. Open Chen, Pan and Chen, Section 2, and write the constitutive law of \eqref{eq:const} against
their numbered equations. Fix the sign convention of $\aME$ and write it in one sentence for the paper.
\item Week 2. Find the published product formula for $\bar\aME$ and replace \eqref{eq:alpha}. Read the
homogenisation section of \cite{Tocci2021} again for the list of constants.
\item Week 3. Read Lim, Zhang and Reddy Section 2 and Barretta and Marotti de Sciarra equations (18) and
(19); confirm the operator placement. Read Quintanilla's equations (2.5) and (3.1)--(3.5).
\item Week 4. Redo Calculations A to D by hand and check the parity statement twice.
\item Week 5. Redo Calculations E to G, build the matrix for the composite at three volume fractions, and
run the ten self-tests.
\item Week 6. Compute the dispersion and attenuation curves and plot the field ratio against the volume
fraction. This is the figure that decides the ridge test at week nine.
\item Week 7 to 8. Write the paper from this file, keeping the worked blocks out of the submitted version.
\item Week 9. Take the first-quartile or second-quartile decision on the ridge test, as agreed.
\end{enumerate}

\bibliographystyle{unsrt}
\bibliography{Lamb_Anu}

\end{document}
